\begin{comment}
% ============================================================
% headers.tex -- intestazioni e piè di pagina (fancyhdr)
% Va incluso dopo \pagestyle{fancyplain}
% ============================================================

\addtolength{\headwidth}{\marginparsep}
\addtolength{\headwidth}{\marginparwidth}
% remember chapter title
\renewcommand{\chaptermark}[1]{\markboth{{\chaptername\ \thechapter}\ \--- #1}{}} %{\markboth{#1}{#1}}
% section number and title
\renewcommand{\sectionmark}[1]{\markright{\thesection\ \ #1}}
\lhead[\fancyplain{\emph{\MioTestoCapitolo}}{\thepage}]{\fancyplain{}{\emph\rightmark}}
\rhead[\fancyplain{\emph{\MioTestoCapitolo}}{\emph\leftmark}]{\fancyplain{}{\thepage}}
\cfoot{}
\end{comment}


% Da qua è fatto da Gemini, e mi fido il giusto
% ============================================================
% headers.tex -- intestazioni e piè di pagina (fancyhdr)
% ============================================================

% COMMENTATO PER EVITARE CHE L'INTESTAZIONE SBORDI DAI BORDI:
% \addtolength{\headwidth}{\marginparsep}
% \addtolength{\headwidth}{\marginparwidth}

\pagestyle{fancy}
\fancyhf{} % Pulisce tutte le intestazioni e piè di pagina

% --- PAGINE NORMALI (Corpo del capitolo) ---
\renewcommand{\chaptermark}[1]{\markboth{{\chaptername\ \thechapter}\ \--- #1}{}}
\renewcommand{\sectionmark}[1]{\markright{\thesection\ \ #1}}

% Intestazioni standard per le pagine interne (es. numero di pagina a sinistra/destra)
\lhead[\fancyplain{}{\thepage}]{\fancyplain{}{\emph\rightmark}}
\rhead[\fancyplain{}{\emph\leftmark}]{\fancyplain{}{\thepage}}
\cfoot{}

% --- PAGINE INIZIALI DEI CAPITOLI (Stile "plain") ---
\fancypagestyle{plain}{
    \fancyhf{} % Pulisce l'intestazione predefinita della prima pagina del capitolo
    
    % Se vuoi il testo a DESTRA sia per pagine pari che dispari (o layout a pagina singola):
    \fancyhead[R]{\emph{\MioTestoCapitolo}}
    
    % (Se invece usi il fronte-retro e lo vuoi sempre in alto a DESTRA nella pagina dispari, usa \fancyhead[RO])
    
    \renewcommand{\headrulewidth}{0pt} % Rimuove la linea orizzontale sotto l'intestazione (opzionale)
}